\section{Définition spectrale de la masse comme mesure de persistance de \(1\) parcours couplé à un régime physique}
La masse est la mesure spectrale de la persistance d'un parcours lorsqu'une section généalogique se couple à un régime physique. Elle n'est pas un attribut scalaire attaché d'emblée à un nom de particule. La dynamique APP--TRIT--TPK construit le parcours, sa mémoire, sa signature et son caractère ; la Mécanique Dimensionnelle réalise cette généalogie dans une droite de masse et détermine, par le couplage, le sous-espace observable et son spectre.

Une particule élémentaire est représentée par une section persistante ; une famille non centrale, par une multisection stable ; une particule composée, par une composition d'enregistrements généalogiques avec frontière de liaison ; une résonance, par un pôle de l'opérateur couplé ; un secteur topologique, par son anneau de fusion et sa représentation de tressage avant toute carte de masse ; une section virtuelle, par son horizon terminal de prolongement. Ces réalisations sont distinctes et ne doivent pas être aplaties en une unique liste de nombres.

La grandeur externe n'intervient jamais dans la sélection de la cellule, du parcours, de la signature, du lecteur, du coefficient de tour ou du vecteur propre. La confrontation métrologique est effectuée après la construction et la fixation du résultat interne.

\Needspace{7\baselineskip}
\section{Espace complet des états de masse}
Soit \(\mathcal C_{81}\) l'atlas cellulaire. Chaque cellule admet quatre orientations two--way, de sorte que l'espace fini des parcours orientés est

\[
 \mathcal R_{324}
 =\{(c,h_0,v_0):c\in\mathcal C_{81},\ (h_0,v_0)\in\{\pm1\}^2\}.
\]
La projection cellulaire, la relation de famille et la stratification grossière produisent

\[
 |\mathcal C_{81}|=81,\qquad
 |\mathcal F_{56}|=56,\qquad
 |\mathcal M_{13}|=13,\qquad
 |\mathcal R_{324}|=324.
\]
L'électron occupe l'unique point fixe central. Les douze multisections non centrales sont des fibres finies stables par inversion cellulaire et n'admettent généralement pas de section ponctuelle équivariante. Leur objet naturel est donc un opérateur sur la fibre, et non un nombre choisi parmi ses cellules. Le secteur nul est incorporé au moyen de la section zéro de la droite de masse, et non comme valeur d'un caractère positif.

Le catalogue physique résultant distingue trois leptons chargés, trois réalisations neutriniques d'interaction, six saveurs de quarks avec leurs orbites de couleur, deux réalisations de jauge nulles, deux bosons de jauge massifs, un secteur scalaire, deux familles composées, trois codomaines topologiques et un système de sections virtuelles. Cette classification ultérieure ne limite pas le nombre d'états composés ou résonants : composition et excitation engendrent des familles supplémentaires au sein de leurs codomaines respectifs.

\Needspace{7\baselineskip}
\section{Trajectoire généalogique d'une réalisation de masse}
La loi ne commence ni par une étiquette de particule ni par un nombre. Son ordre est

\[
\boxed{
\begin{aligned}
\mathrm{APP}\to\mathrm{TRIT}\to\mathrm{TPK}
&\to\widetilde\Gamma^{\rm surv}
\to\Gamma^{\rm obs}
\to L_{108}
\to\sigma\in\mathbb Z^5\\
&\to\chi_0
\to\chi_{\rm full}
\to\widehat{\mathcal M}
\to\text{couplage}\\
&\to\text{spectre/pôle}
\to\mathcal L_M
\to\text{confrontation}.
\end{aligned}}
\]
Ici :

\begin{itemize}
\item APP dispose l'incompatibilité locale, les feuillets somme/produit et leur orientation ;
\item le TRIT conserve signe, seuil et ouverture bidirectionnelle ;
\item le TPK fait évoluer phase, mémoire, retenue, survie et retour ;
\item une particule est une section persistante d'une extension, et son parcours observable est
l'histoire qu'elle laisse dans l'enregistrement généalogique ;

\item la signature entière est additive ;
\item le caractère convertit la composition additive en un rapport positif ;
\item les tours raffinent le caractère au moyen d'un unique opérateur harmonique global ;
\item la droite de masse fournit le type dimensionnel ;
\item mesurer, c'est se coupler : le couplage détermine quel sous-espace et quel spectre sont
physiquement accessibles, sans consulter la grandeur qui sera ensuite mesurée.

\end{itemize}
La masse ne sélectionne pas le parcours. Le parcours persistant, l'enregistrement généalogique et le couplage sélectionnent l'objet spectral dont la masse est lue.

\Needspace{7\baselineskip}
\subsection{Composantes conservées de l'état enrichi du TPK dans la trajectoire de masse}
L'état qui alimente la loi des masses est

\[
 X_K=(w_6,w_{30},R_{36},g_K,B_K,p_K,c_K,\varepsilon_K,\omega_K),
\]
où sont conservés mot court, mot prolongé, région, phase nonadique, mémoire de bloc, préfixe, retenue, feuillet et enroulement. La chaîne

\[
 w_6\longrightarrow w_{30}\longrightarrow R_{36}
 \longrightarrow G_9
\]
ne projette pas directement une masse : elle construit l'identité persistante qui sera ensuite lue.

\Needspace{7\baselineskip}
\subsection{Filtration nonadique avec mémoire}
À partir de chaque état, on obtient une fibre de raffinement formée de \(729\) sous-cylindres,

\[
 \{C_{K,j}:0\le j<729\}.
\]
L'élagage \(G_9\) satisfait

\[
 G_9(C_{K,j})\in\{0,1\},\qquad
 \sum_{j=0}^{728}G_9(C_{K,j})=1.
\]
Si \(j_K\) est le survivant,

\[
 P_{K+1}=729P_K+j_K.
\]
Après neuf pas,

\[
 g(K+9)=g(K),
\]
mais changent, en général,

\[
 (B_{K+9},P_{K+9},\operatorname{pref}_{K+9},
 \operatorname{Witt}_{K+9})
 \ne
 (B_K,P_K,\operatorname{pref}_K,\operatorname{Witt}_K).
\]
Telle est la forme équationnelle de la persistance : la phase revient, et non l'état. La masse peut reconnaître une particule après le retour parce que l'enregistrement généalogique conserve la différence.

\Needspace{7\baselineskip}
\subsection{Transduction de la retenue nonadique vers l'état signé dodécaphasique}
Pour une coordonnée temporelle \(\theta\), nous définissons

\[
 g_0(\theta)=\theta\bmod9,\qquad
 \beta(\theta)=
 \left\lfloor\frac{\widetilde\theta}{9}\right\rfloor\bmod12.
\]
Alors

\[
 g_0(\theta+9)=g_0(\theta),\qquad
 \beta(\theta+9)=\beta(\theta)+1\pmod{12}.
\]
Ainsi, le retour nonadique conserve la phase et avance d'un cran de mémoire. Les 108 temps sont

\[
 108=9\cdot12,
\]
de sorte que chaque enregistrement généalogique contient un cycle complet des deux échelles. Les 36.864 compositions de retenue décrivent les combinaisons locales de cette mémoire bidirectionnelle.

\Needspace{7\baselineskip}
\subsection{Complétion et conservation de l'unité}
La frontière en base mille se décompose comme

\[
 999=729+270=729+243+27,
\]
tandis que la complétion est

\[
 1000=729+271=729+243+27+1,
\qquad
 271=270+1.
\]
Le \(+1\) est l'unique sommet entièrement frontalier et réalise la retenue qui restitue l'unité. La partition \(729/271\) distingue noyau actif et bord médiateur ; elle n'est pas une approximation du nombre d'or. La même conservation permet à un parcours de changer de bloc sans perdre son identité normalisée.

\Needspace{7\baselineskip}
\subsection{Théorème de l'atlas orienté complet}
Chaque cellule \(c\in\mathcal C_{81}\) possède quatre orientations

\[
 (h_0,v_0)\in\{\pm1\}^2.
\]
L’application

\[
 (c,h_0,v_0)\longmapsto\Gamma(c,h_0,v_0)
\]
est bijectif sur \(\mathcal R_{324}\). En conséquence,

\[
 |\mathcal R_{324}|=4|\mathcal C_{81}|=324.
\]
Chaque parcours possède 108 entrées d'enregistrement généalogique, de sorte que le dénombrement temporel total est

\[
 81\cdot108=8748.
\]
Le quadruplet

\[
 (\text{ligne germe},\text{colonne de graines},h_0,v_0)
\]
détermine de manière unique le parcours, son mot \(w_6\), son mot dodécaphasique, ses retours, sa mémoire, sa retenue, sa signature historique et ses lecteurs. Les 324 parcours peuvent donc non seulement être énumérés, mais aussi reconstruits et réunis exactement, champ par champ, sans utiliser de masse.

\Needspace{7\baselineskip}
\subsection{Quotients de famille et de multisection}
Le quotient des 81 cellules par histoire de parcours produit 56 familles. Leurs multiplicités se répartissent comme

\[
 41\times1,\quad10\times2,\quad2\times3,\quad1\times4,\quad2\times5,
\]
et

\[
 41+20+6+4+10=81.
\]
La stratification en treize multisections a pour rangs

\[
 1,\quad2,\quad4,\quad6,\quad8,\quad12,\quad16
\]
avec multiplicités

\[
 1,\quad2,\quad3,\quad2,\quad3,\quad1,\quad1,
\]
et pour somme des rangs

\[
 1+4+12+12+24+12+16=81.
\]
L'électron occupe le bloc de rang un. Les autres secteurs se situent dans des blocs de rang supérieur et exigent un projecteur de représentation, et non une cellule nominale.

\Needspace{7\baselineskip}
\subsection{Signature, lecteurs et mémoire par parcours}
Pour chaque \(\Gamma\in\mathcal R_{324}\), on calcule :

\[
 \sigma_\Gamma\in\mathbb Z^5,\qquad
 K_D(\Gamma),\qquad K_\Omega(\Gamma),
\]
avec les premiers retours de classe, d'orbite, de cellule et d'état cinématique ; la séparation des préfixes ; la dette de retenue ; l'ambiguïté ; les triades stabilisées ; le feuillet et l'enroulement. La famille \(K_D\) présente 14 profils ; \(K_\Omega\), neuf ; ensemble, ils forment 40 paires et 18 coïncidences. Cette diversité constitue l'information fine qui permet de passer d'une fibre grossière à un opérateur physique.

\Needspace{7\baselineskip}
\section{Du défaut APP à la signature entière}
\Needspace{7\baselineskip}
\subsection{Défaut primordial et déploiement}
Au centre APP,

\[
 B_c=9P_++5P_-,\qquad 9-5=4.
\]
Le saut local primordial est 4. Le défaut global \(54\) appartient à un déploiement ultérieur \(2\to6\to54\) ; il ne doit pas être rétroprojeté sur le centre comme s'il s'agissait de la même opération.

\Needspace{7\baselineskip}
\subsection{\(12\) événements et diviseur dodécaphasique}
La coordonnée orientée ne reçoit pas arbitrairement un facteur \(1/6\). La chaîne est

\[
12\ \text{événements}
\to s_C^{\rm tot}
\to 1+5+6
\to s_C^{(12)}=\frac{s_C^{\rm tot}}6
\to W_\pm=6n_A\pm s_C
\to\operatorname{SNF}(1,12)
\to\chi_{\rm mass}.
\]
Le \(1/6\) est le quotient orienté de la clôture de douze événements. Ce n'est ni le résidu (-1/12), ni la pseudo-inverse d'une valence, ni l'incidence hexade--octade.

\Needspace{7\baselineskip}
\subsection{Construction de la signature entière complète à partir de l'enregistrement du parcours}
L'enregistrement généalogique CT--108 produit

\[
 \sigma(\Gamma)
 =\bigl(n_A,s_C,k_{\Delta_4},b_{120},b_\zeta\bigr)\in\mathbb Z^5.
\]
Les coordonnées ont des provenances distinctes :

\begin{itemize}
\item \(n_A\) : durée/persistance dominante ;
\item \(s_C\) : orientation nette de la bicouche, divisée par six après la clôture
dodécaphasique ;

\item \(k_{\Delta_4}\) : canal de torsion minimale ;
\item \(b_{120}\) : coefficient d'événement de l'enregistrement ;
\item \(b_\zeta\) : autocorrélation de pas trois.
\end{itemize}
On définit

\[
 \zeta_{270}:=135\Delta_4^2,
\]
et reste distinct de

\[
 s_{270}=q_-^{270}-q_+^{270}.
\]
Les 324 parcours orientés sont reliés bijectivement à leurs cinq coordonnées historiques par l'identité canonique de parcours. Cette liaison utilise uniquement la structure généalogique et ne consulte aucune grandeur métrologique.

\Needspace{7\baselineskip}
\subsection{Angle direct, contre-angle et canaux conjugués}
La coordonnée angulaire directe procède de la constante de structure fine déjà construite :

\[
 A=1000\alpha_{\rm HMT}
 =7.297352569283800997285105472380663\ldots{}^\circ.
\]
Le contre-angle ne s'obtient ni à partir d'une masse ni à partir de la fonctionnelle de Barbero--Immirzi. Soit

\[
 D_A=\exp\!\left(-\frac{100\pi}{9}\alpha_{\rm HMT}\right),
\]
\[
 \mathcal C_\pi(\alpha_{\rm HMT})
 =169\alpha_{\rm HMT}^{6}
 +\frac{D_A\alpha_{\rm HMT}^{7}}
        {1-D_A\alpha_{\rm HMT}},
\]
\[
 y_*=\frac{\pi}{90}(H_5+\alpha_{\rm HMT})
      -\mathcal C_\pi(\alpha_{\rm HMT}),
 \qquad
 C^*=\frac{180}{\pi}y_*.
\]
Par conséquent,

\[
 C^*=2.123738338968645724261951380393\ldots{}^\circ.
\]
La coordonnée \(A\) est le mode angulaire commun, et \(C^*\) le mode différentiel orienté. Leur passage à la carte archimédienne est

\[
 x=A_{\rm rad}=\frac{\pi A}{180},
 \qquad
 y=C^*_{\rm rad}=\frac{\pi C^*}{180}.
\]
Dans l'algèbre réelle engendrée par l'involution \(R=\operatorname{diag}(1,-1)\), avec \(P_\pm=(I\pm R)/2\), le bloc biangulaire est

\[
 B_{A,C^*}=AI+C^*R
 =(A+C^*)P_++(A-C^*)P_-.
\]
L'angle plus le contre-angle et l'angle moins le contre-angle sont ainsi les deux valeurs propres d'un même objet. L'exponentiation produit les canaux

\[
 q_+=e^{-x-y},\qquad q_-=e^{-x+y}.
\]
La carte est inversible :

\[
 x=-\frac12\log(q_+q_-),
 \qquad
 y=\frac12\log\frac{q_-}{q_+}.
\]
La transformation \(C^*\mapsto-C^*\) conserve les projecteurs et échange les canaux. Après la construction régionale, la coordonnée de retour satisfait

\[
 \eta_{\rm ret}=\frac{C^*}{2}-\alpha_{\rm HMT}\,{\rm deg},
\]
sans constituer une seconde règle de sélection de \(C^*\).

\Needspace{7\baselineskip}
\subsection{Caractère central}
Avec

\[
 \Theta_0=
 \left(A_{\rm rad},\frac{C^*_{\rm rad}}6,
 \Delta_4,s_{120},\zeta_{270}\right),
\]
le caractère central est

\[
 \chi_0(\sigma)=
 \exp\!\left(
 n_AA_{\rm rad}+\frac{s_C}{6}C^*_{\rm rad}
 +k_{\Delta_4}\Delta_4+b_{120}s_{120}+b_\zeta\zeta_{270}
 \right).
\]
Par additivité de la signature,

\[
 \chi_0(\sigma_1+\sigma_2)=\chi_0(\sigma_1)\chi_0(\sigma_2).
\]
L'unité physique n'intervient pas dans cette exponentielle : \(\chi_0\) est un automorphisme positif adimensionnel de la droite de masse.

\Needspace{7\baselineskip}
\section{Hiérarchie harmonique globale des tours}
Avec les coordonnées (x=\allowbreak{}A\_\{\textbackslash{}rm rad\}) et (y=\allowbreak{}C\textasciicircum{}*\_\{\textbackslash{}rm rad\}) déjà construites, soient

\[
 q_\pm=\exp\!\left[-\frac{\pi}{180}(A\pm C^*)\right]
 =e^{-x\mp y}.
\]
Dans une convention globale cohérente,

\[
 T=e^{-xI-yR},\qquad
 c_n=\operatorname{Tr}(T^n)=q_-^n+q_+^n,
\]
\[
 s_n=-\operatorname{Tr}(RT^n)=q_-^n-q_+^n.
\]
De manière équivalente, sous forme opératorielle,

\[
 c_n=2e^{-nx}\cosh(ny),\qquad
 s_n=2e^{-nx}\sinh(ny).
\]
L'involution \(C^*\mapsto-C^*\) fixe les projecteurs et permute les coefficients spectraux ou, de manière équivalente, les canaux associés. On en déduit

\[
 c_{m+n}=\frac{c_mc_n+s_ms_n}{2},\qquad
 s_{m+n}=\frac{s_mc_n+c_ms_n}{2},
\]
\[
 c_n^2-s_n^2=4e^{-2nx},
\]
et la récurrence commune d'ordre deux. Les hauteurs admissibles forment

\[
 \mathfrak S=\langle90,120\rangle
 =\{90,120\}\cup\{30r:r\ge6\}.
\]
Les valeurs \(90,120,180,210,240,270,360,420\) sont des témoins généalogiques, et non la tour complète. Avant le quotient, la hauteur 360 conserve deux histoires :

\[
 k[\mathfrak S_0]
 \simeq k[X_{90},Y_{120}]/(X_{90}^4-Y_{120}^3).
\]
Le relèvement général s'écrit

\[
 \chi_{\rm full}(\sigma,\eta)
 =\chi_0(\sigma)
  \exp\!\left(\sum_{h\in\mathfrak S}\eta_hs_h\right),
\]
avec

\[
 N_{120}=b_{120}+\eta_{120}.
\]
La somme numérique n'efface pas la double provenance : \(b_{120}\) provient de l'enregistrement et \(\eta_{120}\) du raffinement spectral. Elle ne peut pas non plus absorber \(b_\zeta\zeta_{270}\) dans \(\eta_{270}s_{270}\).

\Needspace{7\baselineskip}
\section{Branche absolue : action, horloge et droite de masse}
Le lecteur déterminantal produit

\[
 \operatorname{SNF}(H_4)=\operatorname{diag}(1,2,2,4),\qquad
 |\operatorname{coker}H_4|=16,
\]
\[
 \Sigma_H=\varphi\alpha_{\rm HMT}^{16},\qquad
 \operatorname{ord}_{10}(\Sigma_H)=-34.
\]
Les deux sections orientées d'une seule droite d'action sont

\[
 \hbar_{\rm pre}=H_5\,10^{-34}\mathcal U_S,
 \qquad
 \hbar_{\rm ret}=\eta_{\rm ret}\,10^{-34}\mathcal U_S.
\]
Leur carte additive fine et leur carte multiplicative sont

\[
 \delta_{2w}=H_5-\eta_{\rm ret},\qquad
 R_{\rm act}=\frac{H_5}{\eta_{\rm ret}}>0.
\]
Il ne s'agit pas de deux constantes de Planck. La décennie commune fixe l'ordre absolu ; le quotient transporte l'information relative d'aller/retour.

Une fois \(t_0\in\mathcal L_T^+\) déclaré, CT--108 fixe

\[
 \omega_0=\frac{2\pi}{108t_0},\qquad
 \varepsilon_{\rm clk}=\hbar_{\rm ret}\omega_0,
 \qquad
 m_{\rm clk}=\varepsilon_{\rm clk}c_{\rm int}^{-2}.
\]
Le caractère agit par homothétie sur \(\mathcal L_M\) ; il ne devient pas \(\hbar\). Dans les quotients de masses, tout facteur commun se simplifie :

\[
 \frac{m_Y}{m_X}
 =\chi_{\rm full}(\sigma_Y-\sigma_X,\eta_Y-\eta_X)
 R_{\rm act}^{\kappa_Y-\kappa_X}.
\]
Ainsi, \(10^{-34}\) fixe l'échelle absolue, mais ne corrige pas un résidu relatif. Une modification relative ne peut provenir que du parcours, de l'horloge spécifique, du lecteur ou des tours.

Une distinction reste nécessaire. HMT construit une section interne une fois \(t_0\) déclaré ; la coordonnée MeV exige une carte. Pour déterminer complètement la carte absolue sans convention externe, il faut construire

\[
 \Gamma\longmapsto t_\Gamma
 \quad\text{ou}\quad
 \Gamma\longmapsto\omega_\Gamma.
\]
La section 10 définit un lecteur d'horloge fondé sur les retours et l'holonomie. Sa normalisation absolue est déterminée lorsque l'opérateur de parcours fixe une section positive de la droite temporelle.

\Needspace{7\baselineskip}
